Since the financial markets are very complex and it can be influenced by a natural disaster, the government of a country has been changed, some types of fake news, etc., these kind of interconnected factors matters every single day. Which make forecasting prices very difficult, still it is essential for traders and investment managers. The proposed research tackles these intricacies by combining sophisticated computational approaches including algorithmic trading approaches, technical indicator optimization, deep learning (DL), and deep reinforcement learning (DRL) to forecast the prices of shares on the market, specifically for NVIDIA Corporation (NVDA) and the S\&P500 (GSPC) using available indices data till March 5, 2025.

Concisely stated, a cohort of 20 PineScript-based trading strategies termed "FinDRL" derived from trend-following (TF), breakout detection (BO), mean-reversion (MR), and volume analytics with a stochastic RSI screen, were statistically backtested (NVDA historical price data from 2007–2025). FinDRL - Momentum, FinDRL - Williams \%R, and FinDRL - Volume Spike were among the greatest performers, boasting substantial profits factors (e.g., 2.469) and showcasing the value of algorithmic leverage in complex market conditions. Others were found limiting, emphasizing the need for selected application.

To getting these preliminary results, an optimization framework based on Bayesian optimization (Optuna) was applied to increase the effectiveness of those indicators. Through systematic parameter tuning, the \textbf{IndicatorOptimizer} class has transformed indicators parameter once deemed ineffective, like the Parabolic SAR, into profitable trading strategies, achieving several magnitudes of performance improvement in terms of profit and drawdown reduction. The validation was done in a high fidelity trading environment, making the exercise relevant to real world portfolio management.

Then, more complex DL models were trained and optimized (Bayesian optimization) based on GRU, CNN, and LSTM architectures to improve the forecasting ability even more. These models were incorporated into a dynamic backtesting framework deploying ensemble forecasts, and optimised lookback periods (10–120 days), that consistently surpassed traditional baselines by a wide margin. We used metrics including maximum drawdown and Sharpe ratio, and found significant improvements in predictive accuracy as well as trading capabilities.

The results led to the establishment of a complex DRL system named FinDRL that encapsulated Bayesian-optimized Deep Neural Networks (DNN) used as feature extractors and Deep Q-Networks (DQN) to make actionable decisions within a realistic trading environment (\textbf{FinDRLTradingEnv}). This environment included OHLCV data, Simple Moving Average (SMA) indicators, and transaction costs. Empirical validation showed solid portfolio growth (15–20\%) and good risk-adjusted performance (Sharpe ratio: 1.2). Visualizations Quest visualizations provide comprehensive visualizations (e.g. equity curves, trade signals, drawings, etc.) to allow clear background analysis and comparative performance of strategies.

Overall, this research presents a modular, composable, and scalable architecture that integrates traditional technical indicators with advanced DL and DRL-aligned representations, thereby delivering foundational tools and insights for market participants and researchers.